\section{证据映射总表与上线检查单}\label{sec:mapping}

\subsection{三套规则的证据汇总}

下表把三套规则的关键决策集中列出，按“论文表 5/表 6 直接支持”、“文献支持（论文转引）”、“论文论证”、“实测修正”、“工程判断”五档标注，供上线评审时逐条核对。其中“实测修正”一档来自第~\ref{sec:verify}~节与第~\ref{sec:retune}~节对论文口径的推翻与重标定。

\begin{longtable}{@{}>{\raggedright\arraybackslash}p{2.4cm}>{\raggedright\arraybackslash}p{4.6cm}>{\raggedright\arraybackslash}p{4.4cm}>{\raggedright\arraybackslash}p{2.6cm}@{}}
\caption{三套规则的证据分级总表}\label{tab:allmap}\\
\toprule
\textbf{规则} & \textbf{决策} & \textbf{依据} & \textbf{级别} \\
\midrule
\endfirsthead
\toprule
\textbf{规则} & \textbf{决策} & \textbf{依据} & \textbf{级别} \\
\midrule
\endhead
\bottomrule
\endlastfoot
规则一 & 期货时序用趋势方向 & 论文表 5：全窗口净夏普为正 & 实证直接支持 \\
规则一 & 弃用 $L=5$ & 论文表 5：成本占毛利 $50.1\%$ & 实证直接支持 \\
规则一 & 弃用 $L=250$ & 工程增强后成本占毛利 $108.7\%$（第~\ref{sec:verify}~节） & 实测修正 \\
规则一 & 窗口取 $L=10$（首选）、$20$（备选） & 第~\ref{sec:retune}~节：参数曲面平滑、成本占毛利 $\le18\%$ & 实测修正 \\
规则一 & 波动率目标化 & \cite{barroso2015,harvey2018,moreiramuir2017} & 文献支持 \\
规则一 & 缩放归因检验 & \cite{kim2016,huang2020} & 文献支持 \\
规则一 & 恐慌状态过滤 & \cite{danielmoskowitz2016} & 文献支持 \\
规则一 & 拥挤监控 & \cite{stein2009,loupolk2022,baltas2019} & 文献支持 \\
规则一 & 不做板块中性化 & 论文对偶四（收益来自共同方向风险） & 论文论证 \\
规则一 & 免交易带系数 $0.1$ & — & 工程判断 \\
规则一 & ATR 止损倍数 & — & 工程判断 \\
规则一 & 恐慌阈值 $-5\%$/80 分位 & DM16 定性 + 自标定 & 工程判断 \\
\midrule
规则二 & A 股截面用反转方向 & 论文表 6：$L=20\sim250$ 净夏普为正 & 实证直接支持 \\
规则二 & 取 $L=250$（备选 $60$） & 第~\ref{sec:retune}~节：回撤浅 $16$ 个百分点、成本占毛利低 $6.6$ 个百分点 & 实测修正 \\
规则二 & 月频而非周频 & 论文表 6：$L=5$ 毛利夏普 $1.08$ 被压至 $-0.35$ & 实证直接支持 \\
规则二 & 纯多头 + 指数对冲 & 论文审计章：A 股做空不可实现 & 实证直接支持 \\
规则二 & 行业 + 市值中性化 & \cite{moskowitzgrinblatt1999} & 文献支持 \\
规则二 & 不用波动率缩放 & 论文执行章（机制推理） & 论文论证 \\
规则二 & 单票偏离止损 & \cite{leungli2015} & 文献支持 \\
规则二 & 现金冗余 $30\%$ & \cite{shleifervishny1997,khandanilo2011} & 文献支持 \\
规则二 & 不得逢跌加仓 & \cite{brunnermeier2009} & 文献支持 \\
规则二 & 对冲 beta 窗口长度 & — & 工程判断 \\
规则二 & 行业分类标准 & — & 工程判断 \\
\midrule
规则三 & 期货截面中线用反转 & 论文表 5：$L=60/120$ 净夏普 $0.72/0.56$、$t=2.59/1.98$ & 实证直接支持 \\
规则三 & 弃用 $L\le20$ & 论文表 5：毛利夏普为负或近零 & 实证直接支持 \\
规则三 & 配对的实现方式 & 论文对偶七 + \cite{avellaneda2010} & 论文论证（配对版收益未检验） \\
规则三 & 协整资格检验 & \cite{englegranger1987,dickeyfuller1979} & 文献支持 \\
规则三 & 半衰期筛选 $[5,60]$ & 论文附录 A.1 & 论文论证 \\
规则三 & 带宽 $\propto$ 成本立方根 & \cite{martin2011} & 文献支持 \\
规则三 & Kalman 在线对冲比 & \cite{elliott2005} & 文献支持 \\
规则三 & s-score 进出场阈值 & \cite{avellaneda2010} & 文献支持 \\
规则三 & 止损阈值 $3.5$ & — & 工程判断 \\
规则三 & 产业链划分方案 & — & 工程判断 \\
\end{longtable}

\subsection{仍未补齐的检验}

论文的实证有明确的未竟事项。本文已补做两项——多重检验校正（第~\ref{sec:mtflow}~节）与规则三的配对筛选及回测（第~\ref{sec:pairalgo}~节）；其余四项仍未完成。它们不是可选项，而是规则成立的前提。

\begin{enumerate}[leftmargin=2.2em]
  \item \textbf{涨跌停与停牌的可交易性建模}。论文未做。需要在信号与收益计算中显式排除涨跌停板与停牌日的成交，否则 2015 年等异常波动时段的可实现收益被系统性高估。本文的规则文本已写入顺延规则，但回测引擎尚未接入（第~\ref{sec:mtflow}~节第 7 步因此判为未通过）。
  \item \textbf{参与率约束的冲击成本}。论文用固定基点。反转类规则（规则二、三）的换手高且交易密集于流动性差的时点，必须改用 \mbox{impact}$=K\sqrt{\text{参与率}}$ 模型并标定 $K$。依据：\cite{almgrenchriss2000,perold1988}。
  \item \textbf{A 股复权因子的补全}。论文用“跌幅超涨跌停即判除权”的近似规则，漏掉小额分红型除权。需从数据商处取得复权因子序列替换。
  \item \textbf{规则一的缩放与信号分离归因}。分别测试“只保留符号”与“只做缩放”的对照组，确认收益来源。依据：\cite{kim2016,huang2020}。
\end{enumerate}

\subsection{持续监控清单}

规则上线后需持续监控的四类指标，以及触发时的动作：

\begin{table}[htbp]
\centering\footnotesize
\caption{持续监控清单}\label{tab:monitor}
\begin{tabular}{@{}>{\raggedright\arraybackslash}p{3.4cm}>{\raggedright\arraybackslash}p{4.6cm}>{\raggedright\arraybackslash}p{3.6cm}>{\raggedright\arraybackslash}p{2.4cm}@{}}
\toprule
\textbf{监控项} & \textbf{指标定义} & \textbf{触发阈值} & \textbf{动作} \\
\midrule
成本占比 & 实测成本占毛利比 & 规则一 $>35\%$；规则二 $>40\%$ & 切换更长窗口或停机 \\
拥挤度 & 策略内两两平均相关（滚动 3 年分位） & $>90$ 分位 & 规则一敞口 $\times0.7$ \\
恐慌状态 & 市场 60 日收益与波动率分位 & 收益 $<-5\%$ 且波动 $>$80 分位 & 规则一敞口 $\times0.5$ \\
实测偏度 & 日收益偏度 & 趋势规则转负且显著 & 复核是否已失去凸性 \\
最大回撤 & 组合净值回撤 & 规则一 $15\%/25\%$；规则二三 $20\%/30\%$ & 腰斩 / 全平 \\
资金冗余 & 现金或保证金缓冲比例 & 低于设计值 & 降低敞口而非补足杠杆 \\
行业偏离 & 实际行业权重与基准偏离 & 单行业 $>5\%$ & 调整持仓 \\
产业链集中度 & 单产业链持仓占比 & $>50\%$ & 调整配对分布 \\
\bottomrule
\end{tabular}
\end{table}

\subsection{最后一条提醒}

本文给出的三套规则都有一个共同的性质：\textbf{它们的历史绩效来自一个未做可交易性过滤、成本模型简化的对照实验}。论文已经把这些局限逐条披露（审计章），本文在规则设计层面对其中三项做了补救——可交易性处理、成本门槛、失效防线，并在第~\ref{sec:mtflow}~节补做了多重检验校正。但回到可实现性，第~\ref{sec:mapping}~节列出的四项仍未完成：涨跌停与停牌过滤尚未接入回测引擎、冲击成本模型未标定、A 股复权因子未补全、规则一的缩放归因未分离。

在补齐之前，这三套规则应当被当作\emph{设计模板}而不是\emph{可上线策略}。把模板误当作策略，正是论文反复警告的那类错误——用统计上尚不足以称为显著的结果，去支撑实际的资金决策。
